\documentclass[tikz, border=2pt]{standalone}

\usepackage{amsmath,amssymb}
\usepackage{pgfplots}
\pgfplotsset{compat=1.18}
\usetikzlibrary{arrows.meta}
\usepackage{xcolor}

\definecolor{utnblue}{RGB}{0, 51, 102}

\begin{document}
\begin{tikzpicture}[>=Latex, font=\small]
\begin{axis}[
    axis lines=middle,
    axis line style={-{Latex[length=4pt]}},
    width=9cm, height=4.2cm,
    xmin=-5.5, xmax=0.7,
    ymin=-0.75, ymax=0.85,
    xlabel={$\sigma$}, ylabel={$j\omega$},
    xlabel style={anchor=north west},
    ylabel style={anchor=south west},
    xtick={-3,-2,-1}, xticklabels={$-3$,$-2$,$-1$},
    ytick=\empty,
    tick style={black},
    xticklabel style={font=\footnotesize, yshift=-1pt},
    clip=false,
]
    % --- ramas del lugar (sobre el eje real) con sentido de k creciente ---
    % rama desde el polo -1 hacia el cero en -2
    \draw[-{Latex[length=5pt]}, line width=1.4pt, red!80!black]
        (axis cs:-1,0) -- (axis cs:-1.92,0);
    % rama desde el polo -3 hacia -infinito
    \draw[-{Latex[length=5pt]}, line width=1.4pt, red!80!black]
        (axis cs:-3,0) -- (axis cs:-5.25,0);

    % --- polos (x) y cero (o), encima de las ramas ---
    \addplot[only marks, mark=x, mark size=4pt, line width=1pt, utnblue]
        coordinates {(-1,0) (-3,0)};
    \addplot[only marks, mark=o, mark size=4pt, line width=1pt, utnblue]
        coordinates {(-2,0)};

    % --- k=0 en los polos de arranque ---
    \node[anchor=south, font=\scriptsize] at (axis cs:-1,0.1) {$k=0$};
    \node[anchor=south, font=\scriptsize] at (axis cs:-3,0.1) {$k=0$};
\end{axis}
\end{tikzpicture}
\end{document}
